\documentclass[11pt,aspectratio=169]{beamer}

% Theme settings
\usetheme[progressbar=frametitle, block=fill]{metropolis}
\useoutertheme{metropolis}
\useinnertheme{metropolis}
\usecolortheme{metropolis}
\setbeamercolor{background canvas}{bg=white}
\setbeamertemplate{navigation symbols}{}

\usepackage{graphicx}
\usepackage{booktabs}
\usepackage{amsmath,amssymb}

% Define relative path for figures assuming compilation from starter_pack/slides/
\graphicspath{{../figures/}}

\title[From Linear Scores to a Single Hidden Layer]{From Linear Scores to a Single Hidden Layer:\\A Mathematical Study of Simple Learning Systems}
\subtitle{Math4AI Final Capstone}
\author[Team]{Saida Arabova \and Laman Mirzayeva \and Nigar Rustamova  \and Nazrin Aliyeva}
\institute{National AI Center - AI Academy}
\date{March 30, 2026}

\begin{document}

% Slide 1: Title
\begin{frame}
    \titlepage
\end{frame}

% Slide 2: The Central Question
\begin{frame}{The Central Question}
    \begin{block}{Our Objective}
        \textit{When does a one-hidden-layer nonlinear classifier genuinely improve on a linear decision rule, and when is additional model complexity unnecessary?}
    \end{block}
    
    \vspace{0.5em}
    \textbf{Models Compared:}
    \begin{itemize}
        \item \textbf{Softmax Regression (SR):} Linear score function, affine decision boundaries.
        \item \textbf{1-Hidden-Layer Neural Network (NN):} Non-linear representation via $\tanh$ activation, followed by linear classification.
    \end{itemize}
    
    \vspace{0.5em}
    All components (forward pass, backpropagation, and optimization) were implemented entirely from scratch in NumPy.
\end{frame}

% Slide 3: Ensuring a Fair Comparison
\begin{frame}{Ensuring a Fair Comparison}
    Because the models differ in structure and parameterization, our comparison must be mathematically fair.
    
    \vspace{1em}
    \textbf{Our Experimental Contract:}
    \begin{itemize}
        \item \textbf{Identical Splits:} Fixed Train / Validation / Test indices across all runs.
        \item \textbf{Identical Preprocessing:} Digits used exactly as provided (scaled to $[0,1]$), with no task-specific feature engineering.
        \item \textbf{Strict Model Selection:} Settings and checkpoints chosen \textit{only} based on Validation Cross-Entropy.
        \item \textbf{Evaluation Metrics:} Unregularized Test Cross-Entropy (measuring predictive distribution quality) and Test Accuracy.
        \item \textbf{Statistical Rigor:} Final evaluations on the digits benchmark use 5 repeated random seeds.
    \end{itemize}
\end{frame}

% Slide 4: When is Linear Structure Sufficient?
\begin{frame}{When is Linear Structure Sufficient?}
    \textbf{Task:} Linear Gaussian (2 classes, 2D)
    
    \begin{columns}
        \begin{column}{0.5\textwidth}
            \centering
            \includegraphics[height=0.6\textheight]{figures/gaussian_sr.png}\\
            \small Softmax Regression
        \end{column}
        \begin{column}{0.5\textwidth}
            \centering
            \includegraphics[height=0.6\textheight]{figures/gaussian_nn.png}\\
            \small Neural Network ($m=32$)
        \end{column}
    \end{columns}
    
    \vspace{0.5em}
    \textbf{Conclusion:} When the true Bayes-optimal separator is a hyperplane, adding a hidden layer provides \textbf{no meaningful advantage}.
\end{frame}

% Slide 5: Where Nonlinearity Helps
\begin{frame}{Where Nonlinearity Helps}
    \textbf{Task:} Moons (2 classes, 2D interlocking crescents)
    
    \begin{columns}
        \begin{column}{0.5\textwidth}
            \centering
            \includegraphics[height=0.6\textheight]{figures/moons_sr.png}\\
            \small Softmax Regression (Fails)
        \end{column}
        \begin{column}{0.5\textwidth}
            \centering
            \includegraphics[height=0.6\textheight]{figures/moons_nn_optimized.png}\\
            \small Neural Network (Succeeds)
        \end{column}
    \end{columns}
    
    \vspace{0.5em}
    \textbf{Conclusion:} The NN's hidden layer maps the 2D input into a 32D space where the classes can be linearly separated, solving a problem SR intrinsically cannot.
\end{frame}

% Slide 6: Results on Digits Benchmark
\begin{frame}{Results on Digits Benchmark}
    \textbf{Task:} Classify $8 \times 8$ pixel digits (64 dimensions, 10 classes).
    
    \vspace{1em}
    \begin{table}
        \centering
        \begin{tabular}{lcccc}
            \toprule
            & \multicolumn{2}{c}{\textbf{Softmax Regression}} &
              \multicolumn{2}{c}{\textbf{1-Hidden-Layer NN}} \\
            \cmidrule(lr){2-3}\cmidrule(lr){4-5}
            \textbf{Split (Seed 0)} & Accuracy & CE & Accuracy & CE \\
            \midrule
            Train      & 0.965 & 0.202 & 0.991 & 0.063 \\
            Validation & 0.946 & 0.214 & 0.972 & 0.103 \\
            Test       & 0.938 & 0.268 & 0.951 & 0.162 \\
            \bottomrule
        \end{tabular}
    \end{table}
    
    \vspace{0.5em}
    While Softmax Regression performs reasonably well, the Neural Network fits the training distribution more precisely and achieves superior generalization.
\end{frame}

% Slide 7: What Repeated Seeds Add
\begin{frame}{What Repeated-Seed Evaluation Adds}
    A single run can be lucky or unlucky due to random initialization. Re-running with 5 seeds provides statistical confidence.
    
    \vspace{1em}
    \begin{table}
        \centering
        \begin{tabular}{lcc}
            \toprule
            \textbf{Model} & \textbf{Test Accuracy (95\% CI)} & \textbf{Test CE (95\% CI)} \\
            \midrule
            Softmax Regression  & $0.9386 \pm 0.0030$ & $0.2684 \pm 0.0027$ \\
            1-Hidden-Layer NN   & $0.9538 \pm 0.0024$ & $0.1498 \pm 0.0063$ \\
            \bottomrule
        \end{tabular}
    \end{table}
    
    \vspace{1em}
    \textbf{Takeaways:}
    \begin{itemize}
        \item Non-overlapping confidence intervals confirm the NN's advantage is statistically robust.
        \item SR is highly stable but bounded by its linear representational ceiling.
        \item The NN has higher variance in its predictive distributions (Cross-Entropy).
    \end{itemize}
\end{frame}

% Slide 8: Failure Case Analysis
\begin{frame}{Failure Case Analysis: Capacity Limits}
    \textbf{Experiment:} Moons task with a severely restricted hidden width ($m=2$).
    
    \begin{columns}
        \begin{column}{0.4\textwidth}
            \centering
            \includegraphics[width=\textwidth]{figures/ablation_h2.png}
        \end{column}
        \begin{column}{0.6\textwidth}
            \begin{itemize}
                \item Even with 1000 epochs (5x budget) and an optimized learning rate, the boundary remains near-linear.
                \item \textbf{Why?} A 2-dimensional nonlinear manifold cannot solve a non-convex partition of 2D input space.
                \item \textbf{Lesson:} This is an \textit{under-capacity} failure, not an optimization failure. No amount of training time can overcome a structural representational limit.
            \end{itemize}
        \end{column}
    \end{columns}
\end{frame}

% Slide 9: Track B - Prediction Confidence
\begin{frame}{Track B: Prediction Confidence and Reliability}
    \textit{Does the model know when it is likely to be wrong?}
    
    \begin{columns}
        \begin{column}{0.4\textwidth}
            \centering
            \includegraphics[width=\textwidth]{figures/track_b_reliability_diagram.png}
        \end{column}
        \begin{column}{0.6\textwidth}
            \begin{itemize}
                \item \textbf{Softmax Regression} is systematically underconfident (empirical accuracy exceeds stated probability).
                \item \textbf{The Neural Network} is well-calibrated. In its highest confidence bin $[0.80, 1.00]$, its mean confidence is 97.5\% and empirical accuracy is 99.1\% (Gap = $-0.016$).
                \item The NN encodes uncertainty more efficiently: sharply confident when right, higher entropy when wrong.
            \end{itemize}
        \end{column}
    \end{columns}
\end{frame}

% Slide 10: Limitations
\begin{frame}{Limitations of Our Study}
    \begin{itemize}
        \item \textbf{Dataset Scale:} The digits dataset contains only 1,797 samples. Our findings may not seamlessly extrapolate to large-scale vision or language regimes.
        \item \textbf{Architectural Depth:} We analyzed exactly \textit{one} hidden layer. We make no claims about the optimization dynamics or representational power of deeper networks.
        \item \textbf{Low-Dimensional Geometry:} The clear visual intuitions from the 2D Gaussian and Moons tasks do not perfectly mirror the geometric realities of high-dimensional feature spaces.
    \end{itemize}
\end{frame}

% Slide 11: Summary
\begin{frame}{Summary}
    \begin{itemize}
        \item \textbf{When is linear sufficient?} When the data geometry dictates an affine separator (e.g., Gaussian task).
        \item \textbf{What does the hidden layer add?} It warps the input space, allowing a linear output layer to solve nonlinearly separable problems (e.g., Moons, Digits).
        \item \textbf{Is complexity always justified?} No. If structural capacity is too low (our failure case), or the task geometry is too simple, the extra complexity adds nothing.
        \item \textbf{Overall:} The One-Hidden-Layer NN outperformed Softmax Regression significantly on the digits benchmark in both raw accuracy and calibrated reliability.
    \end{itemize}
    
    \vspace{1em}
    \begin{center}
        \textbf{Thank You! \\ Questions?}
    \end{center}
\end{frame}

\end{document}
